\documentclass[11pt]{article}
\usepackage[a4paper,margin=1in]{geometry}
\usepackage{amsmath,amssymb,amsthm,mathtools}
\usepackage[T1]{fontenc}
\usepackage{lmodern}
\usepackage{microtype}
\usepackage{enumitem}
\usepackage{hyperref}

\hypersetup{colorlinks=true,linkcolor=blue,citecolor=blue,urlcolor=blue}

\newtheorem{theorem}{Theorem}[section]
\newtheorem{lemma}[theorem]{Lemma}
\newtheorem{corollary}[theorem]{Corollary}
\newtheorem{proposition}[theorem]{Proposition}
\newtheorem{remark}[theorem]{Remark}
\newtheorem{definition}[theorem]{Definition}

\newcommand{\Z}{\mathbf Z}
\newcommand{\Q}{\mathbf Q}
\newcommand{\Fp}{\mathbf F_p}
\newcommand{\Zp}{\mathbf Z_{(p)}}
\newcommand{\ord}{\operatorname{ord}}
\newcommand{\Bern}{B}
\newcommand{\floor}[1]{\left\lfloor #1\right\rfloor}

\title{An Elementary Voronoi Proof of the Unrestricted Kummer Congruence for Bernoulli Numbers}
\author{}
\date{}

\begin{document}
\maketitle

\section*{Purpose}

This note proves the Bernoulli congruence needed in the irregular-prime
argument, without using $p$-adic $L$-functions, class field theory,
Chebotarev, Hilbert symbols, or cyclotomic units.  The proof is the classical
Voronoi proof of Kummer's congruence, but written carefully enough to remove
side conditions such as $p\nmid m$ or $p\nmid m+1$.

The main theorem is the following.

\begin{theorem}[Kummer congruence for $\Bern_m/m$]\label{thm:kummer}
Let $p\ge 5$ be prime.  Let $m,n>0$ be even integers.  Assume
\[
  m\equiv n \pmod {p-1}
  \qquad\text{and}\qquad
  (p-1)\nmid n.
\]
Then $\Bern_m/m$ and $\Bern_n/n$ are $p$-integral and
\[
  \frac{\Bern_m}{m}\equiv \frac{\Bern_n}{n}\pmod p.
\]
Equivalently,
\[
  \frac{\Bern_m}{m}-\frac{\Bern_n}{n}\in p\Z_{(p)}.
\]
\end{theorem}

The important point is that no hypothesis like $p\nmid m$, $p\nmid n$,
$p\nmid m+1$, or $p\nmid n+1$ is present.  This is exactly what is needed when
large indices such as $m=p r+1$ or $m=r+1$ may be divisible by $p$.

We also extract the special comparison that often appears as Diekmann's
Corollary 34.

\begin{corollary}[The comparison $\Bern_{pr+1}$ versus $\Bern_{r+1}/(r+1)$]
\label{cor:pr-plus-one}
Let $p\ge 5$ be prime.  Let $r>0$ be odd and assume
\[
  (p-1)\nmid r+1.
\]
Then
\[
  \Bern_{pr+1}\equiv \frac{\Bern_{r+1}}{r+1}\pmod p.
\]
Equivalently,
\[
  \Bern_{pr+1}-\frac{\Bern_{r+1}}{r+1}\in p\Z_{(p)}.
\]
\end{corollary}

\section{Conventions and basic inputs}

We use the Bernoulli numbers defined by
\[
  \frac{T}{e^T-1}=\sum_{n\ge 0}\Bern_n\frac{T^n}{n!},
\]
so that
\[
  \Bern_1=-\frac12.
\]

For rational numbers $x,y\in\Q$, we write
\[
  x\equiv y\pmod p
\]
if
\[
  x-y\in p\Z_{(p)},
\]
where
\[
  \Z_{(p)}=\left\{\frac ab\in\Q: p\nmid b\right\}
\]
is the localization of $\Z$ at $p$.

We use two standard elementary facts about Bernoulli numbers.

\begin{theorem}[Faulhaber's formula]\label{thm:faulhaber}
For $h\ge 1$ and $N\ge 1$,
\[
  \sum_{x=1}^{N-1}x^h
  =
  \frac{1}{h+1}
  \sum_{j=0}^{h}
  {h+1\choose j}\Bern_j N^{h+1-j}.
\]
\end{theorem}

\begin{theorem}[von Staudt--Clausen denominator consequence]\label{thm:vsc}
For every $j\ge 0$ and every prime $p$, one has
\[
  v_p(\Bern_j)\ge -1.
\]
Moreover, if $j>0$ is even, then
\[
  v_p(\Bern_j)=-1
  \quad\Longleftrightarrow\quad
  (p-1)\mid j.
\]
In particular, if $j>0$ is even and $(p-1)\nmid j$, then
\[
  \Bern_j\in \Z_{(p)}.
\]
\end{theorem}

\begin{remark}
For the proof below, we only need the displayed consequences of
von Staudt--Clausen.  If a formal development already has
von Staudt--Clausen, this theorem is the only Bernoulli-denominator input.
\end{remark}

\section{A strong power-sum congruence}

The key technical point is that the usual congruence
\[
  \sum_{x=1}^{p-1}x^h\equiv p\Bern_h\pmod {p^2}
\]
is not strong enough when $p\mid h$.  To divide by $h$ later in Voronoi's
argument, we need the stronger modulus $h p^2$.

\begin{lemma}[A valuation estimate]\label{lem:valuation-estimate}
Let $p\ge 5$ be prime and let $s\ge 2$.  Then
\[
  s-v_p(s)-v_p(s+1)\ge 2.
\]
\end{lemma}

\begin{proof}
The cases $s=2,3,4$ are immediate, since $p\ge5$.
Assume $s\ge5$.  For any integer $N\ge4$, we have
\[
  v_p(N)\le \frac{N}{4}.
\]
Indeed, if $v_p(N)=0$ this is clear.  If $v_p(N)=e\ge1$, then
$p^e\mid N$, so $5^e\le p^e\le N$, and $4e\le 5^e\le N$.
Therefore
\[
  v_p(s)+v_p(s+1)
  \le \frac{s}{4}+\frac{s+1}{4}
  =\frac{2s+1}{4}.
\]
For $s\ge5$,
\[
  \frac{2s+1}{4}\le s-2.
\]
Hence
\[
  v_p(s)+v_p(s+1)\le s-2,
\]
which is equivalent to the desired inequality.
\end{proof}

\begin{lemma}[Strong power-sum congruence]\label{lem:power-sum-strong}
Let $p\ge5$ be prime.  Let $h>0$ be even and assume
\[
  (p-1)\nmid h.
\]
Put
\[
  S_h(p)=\sum_{x=1}^{p-1}x^h.
\]
Then
\[
  S_h(p)-p\Bern_h\in h p^2\Z_{(p)}.
\]
Equivalently,
\[
  S_h(p)\equiv p\Bern_h\pmod {h p^2}.
\]
\end{lemma}

\begin{proof}
By Faulhaber's formula with $N=p$,
\[
  S_h(p)
  =
  \frac{1}{h+1}
  \sum_{j=0}^{h}{h+1\choose j}\Bern_j p^{h+1-j}.
\]
The term with $j=h$ is exactly
\[
  p\Bern_h.
\]
It remains to show that every term with $j<h$ lies in $h p^2\Z_{(p)}$.

Write $j=h-s$, where $1\le s\le h$.  Then the corresponding term is
\[
  T_s
  =
  {h\choose s}\frac{\Bern_{h-s}}{s+1}p^{s+1},
\]
because
\[
  \frac{1}{h+1}{h+1\choose h-s}
  =
  \frac{1}{s+1}{h\choose s}.
\]

First consider $s=1$.  Then
\[
  T_1=\frac{h}{2}\Bern_{h-1}p^2.
\]
If $h>2$, then $h-1>1$ is odd, so $\Bern_{h-1}=0$.
If $h=2$, then $\Bern_{1}=-1/2$, and
\[
  T_1=-\frac{h p^2}{4}\in h p^2\Z_{(p)},
\]
since $p\ge5$.

Now let $s\ge2$.  By von Staudt--Clausen,
\[
  v_p(\Bern_{h-s})\ge -1.
\]
Also
\[
  {h\choose s}=\frac{h}{s}{h-1\choose s-1},
\]
so
\[
  v_p\left({h\choose s}\right)
  \ge v_p(h)-v_p(s).
\]
Therefore
\[
\begin{aligned}
  v_p(T_s)
  &\ge
  v_p(h)-v_p(s)+(s+1)-v_p(s+1)-1 \\
  &=
  v_p(h)+s-v_p(s)-v_p(s+1).
\end{aligned}
\]
By Lemma~\ref{lem:valuation-estimate},
\[
  s-v_p(s)-v_p(s+1)\ge2.
\]
Hence
\[
  v_p(T_s)\ge v_p(h)+2,
\]
which is precisely
\[
  T_s\in h p^2\Z_{(p)}.
\]
Thus every term except $p\Bern_h$ lies in $hp^2\Z_{(p)}$, and the lemma follows.
\end{proof}

\section{Voronoi's congruence without auxiliary side conditions}

We now prove the congruence that removes the unwanted restrictions such as
$p\nmid h$ or $p\nmid h+1$.

\begin{theorem}[Voronoi congruence]\label{thm:voronoi}
Let $p\ge5$ be prime.  Let $h>0$ be even and assume
\[
  (p-1)\nmid h.
\]
Let $a\in\Z$ be prime to $p$.  For $1\le x\le p-1$, set
\[
  q_x=\floor{\frac{a x}{p}}.
\]
Then
\[
  \boxed{
  (a^h-1)\frac{\Bern_h}{h}
  \equiv
  a^{h-1}\sum_{x=1}^{p-1}q_x x^{h-1}
  \pmod p.}
\]
In particular, if $a^h\not\equiv1\pmod p$, then $\Bern_h/h\in\Z_{(p)}$.
\end{theorem}

\begin{proof}
For $1\le x\le p-1$, define the least positive residue
\[
  r_x=a x-pq_x.
\]
Since $a$ is prime to $p$, the numbers $r_x$ are a permutation of
$1,2,\dots,p-1$.  Hence
\[
  \sum_{x=1}^{p-1}r_x^h=S_h(p).
\]

We claim that, for each $x$,
\[
  r_x^h
  \equiv
  a^h x^h-hp a^{h-1}q_x x^{h-1}
  \pmod {h p^2\Z_{(p)}}.
\]
Indeed, expanding
\[
  r_x^h=(a x-pq_x)^h
\]
by the binomial theorem, the terms of order $0$ and $1$ in $pq_x$ are exactly
\[
  a^h x^h
  \quad\text{and}\quad
  -hp a^{h-1}q_x x^{h-1}.
\]
For a term of order $\nu\ge2$, the coefficient contains
\[
  {h\choose \nu}p^\nu.
\]
Since
\[
  {h\choose \nu}=\frac{h}{\nu}{h-1\choose \nu-1},
\]
we have
\[
  v_p\left({h\choose \nu}p^\nu\right)
  \ge v_p(h)-v_p(\nu)+\nu.
\]
For $\nu\ge2$ and $p\ge5$, one has
\[
  \nu-v_p(\nu)\ge2.
\]
Therefore every higher-order term is in $h p^2\Z_{(p)}$.

Summing over $x$ gives
\[
  S_h(p)
  \equiv
  a^h S_h(p)-hp a^{h-1}\sum_{x=1}^{p-1}q_xx^{h-1}
  \pmod {h p^2\Z_{(p)}}.
\]
Thus
\[
  (a^h-1)S_h(p)
  \equiv
  hp a^{h-1}\sum_{x=1}^{p-1}q_xx^{h-1}
  \pmod {h p^2\Z_{(p)}}.
\]
By Lemma~\ref{lem:power-sum-strong},
\[
  S_h(p)\equiv p\Bern_h\pmod {h p^2\Z_{(p)}}.
\]
Substituting this congruence yields
\[
  (a^h-1)p\Bern_h
  \equiv
  hp a^{h-1}\sum_{x=1}^{p-1}q_xx^{h-1}
  \pmod {h p^2\Z_{(p)}}.
\]
Dividing by $hp$ in $\Q$ gives
\[
  (a^h-1)\frac{\Bern_h}{h}
  \equiv
  a^{h-1}\sum_{x=1}^{p-1}q_xx^{h-1}
  \pmod p.
\]
This proves the congruence.

If $a^h\not\equiv1\pmod p$, then $a^h-1$ is a unit in $\Z_{(p)}$.
The right-hand side is an integer, so the congruence implies that
$\Bern_h/h\in\Z_{(p)}$.
\end{proof}

\section{Kummer's congruence}

\begin{theorem}[Unrestricted Kummer congruence]\label{thm:kummer-proof}
Let $p\ge5$ be prime.  Let $m,n>0$ be even integers.  Assume
\[
  m\equiv n\pmod {p-1}
  \qquad\text{and}\qquad
  (p-1)\nmid n.
\]
Then
\[
  \frac{\Bern_m}{m},\frac{\Bern_n}{n}\in\Z_{(p)}
\]
and
\[
  \frac{\Bern_m}{m}\equiv\frac{\Bern_n}{n}\pmod p.
\]
\end{theorem}

\begin{proof}
Because $m\equiv n\pmod {p-1}$, the condition $(p-1)\nmid n$ also implies
\[
  (p-1)\nmid m.
\]
Choose an integer $a$ whose image in $\Fp^\times$ is a generator.  Then
\[
  a^m\not\equiv1\pmod p
  \qquad\text{and}\qquad
  a^n\not\equiv1\pmod p.
\]
By Voronoi's congruence,
\[
  (a^m-1)\frac{\Bern_m}{m}
  \equiv
  a^{m-1}\sum_{x=1}^{p-1}\floor{\frac{a x}{p}}x^{m-1}
  \pmod p,
\]
and
\[
  (a^n-1)\frac{\Bern_n}{n}
  \equiv
  a^{n-1}\sum_{x=1}^{p-1}\floor{\frac{a x}{p}}x^{n-1}
  \pmod p.
\]
Since
\[
  m\equiv n\pmod {p-1},
\]
we have
\[
  a^m\equiv a^n\pmod p,
  \qquad
  a^{m-1}\equiv a^{n-1}\pmod p,
\]
and for every $1\le x\le p-1$,
\[
  x^{m-1}\equiv x^{n-1}\pmod p.
\]
Thus the right-hand sides of the two Voronoi congruences are equal modulo
$p$, and the factors $a^m-1$ and $a^n-1$ are equal nonzero elements of
$\Fp$.  Therefore
\[
  \frac{\Bern_m}{m}\equiv\frac{\Bern_n}{n}\pmod p.
\]
The $p$-integrality of both quotients follows from the final assertion of
Theorem~\ref{thm:voronoi}, since $a^m-1$ and $a^n-1$ are units modulo $p$.
\end{proof}

This is exactly Theorem~\ref{thm:kummer}.

\section{The special comparison $\Bern_{pr+1}\equiv\Bern_{r+1}/(r+1)$}

We now prove the exact corollary needed in the infinitude argument.

\begin{corollary}\label{cor:special}
Let $p\ge5$ be prime.  Let $r>0$ be odd and assume
\[
  (p-1)\nmid r+1.
\]
Then
\[
  \Bern_{pr+1}\equiv \frac{\Bern_{r+1}}{r+1}\pmod p.
\]
\end{corollary}

\begin{proof}
Set
\[
  m=pr+1,
  \qquad
  n=r+1.
\]
Since $p$ and $r$ are both odd, $pr$ is odd, so $m=pr+1$ is even.
Also $n=r+1$ is even.  Moreover,
\[
  m=pr+1\equiv r+1=n\pmod {p-1},
\]
because $p\equiv1\pmod {p-1}$.  By hypothesis,
\[
  (p-1)\nmid n.
\]
Therefore Theorem~\ref{thm:kummer-proof} gives
\[
  \frac{\Bern_{pr+1}}{pr+1}
  \equiv
  \frac{\Bern_{r+1}}{r+1}
  \pmod p.
\]
But
\[
  pr+1\equiv1\pmod p,
\]
so multiplication by $(pr+1)^{-1}$ is the same as multiplication by $1$
modulo $p$.  Hence
\[
  \frac{\Bern_{pr+1}}{pr+1}
  \equiv
  \Bern_{pr+1}
  \pmod p.
\]
Combining these congruences gives
\[
  \Bern_{pr+1}
  \equiv
  \frac{\Bern_{r+1}}{r+1}
  \pmod p.
\]
\end{proof}

\section{Lean-facing theorem statements}

The main theorem should be formalized as an unrestricted Kummer congruence,
not as a Voronoi theorem with auxiliary side conditions.

A possible target statement is:
\[
\begin{gathered}
  p\ge5,\quad p\text{ prime},\quad m,n>0,\quad m,n\text{ even},\\
  m\equiv n\pmod {p-1},\quad (p-1)\nmid n
\end{gathered}
\]
implies
\[
  \exists z\in\Z_{(p)},\qquad
  \frac{\Bern_m}{m}-\frac{\Bern_n}{n}=p z.
\]

The special corollary should be stated as:
\[
\begin{gathered}
  p\ge5,\quad p\text{ prime},\quad r>0,\quad r\text{ odd},\quad (p-1)\nmid r+1
\end{gathered}
\]
implies
\[
  \exists z\in\Z_{(p)},\qquad
  \Bern_{pr+1}-\frac{\Bern_{r+1}}{r+1}=p z.
\]

The dependency graph is:
\[
\boxed{\text{Faulhaber}}
\quad+
\boxed{\text{von Staudt--Clausen denominator bound}}
\Longrightarrow
\boxed{S_h(p)\equiv p\Bern_h\pmod {hp^2}}
\]
\[
\Longrightarrow
\boxed{\text{Voronoi congruence without side conditions}}
\Longrightarrow
\boxed{\text{Kummer congruence}}
\Longrightarrow
\boxed{\Bern_{pr+1}\equiv\Bern_{r+1}/(r+1)\pmod p}.
\]

This proof is elementary: it uses only Bernoulli power-sum formulas,
von Staudt--Clausen, the binomial theorem, and finite-field cyclicity.
It does not use $p$-adic $L$-functions, class field theory, Chebotarev,
cyclotomic units, Hilbert symbols, or reciprocity laws.

\end{document}
